\icode{E-T-CA-L2-260625-001}
小汽车正常行驶时，胎压需要稳定在\,\qty{220}{kPa}至\,\qty{280}{kPa}之间。在冬季，某室内停车场温度为\,\qty{27}{\degreeCelsius}，此时汽车仪表盘显示左前轮胎压为\,\qty{252}{kPa}，若将轮胎内气体视为理想气体，热力学温度\(T\)与摄氏温度\(t\)之间的数量关系为\(T=t +\)\qty{273}{K}，忽略轮胎体积的变化，求：
\begin{enumerate}[(1)]
    \item 若室外温度为\,\qty{-23}{\degreeCelsius}，司机将车停在室外足够长时间后，通过计算说明胎压是否符合正常行驶要求（假设轮胎不漏气）；
    \item 汽车行驶一段时间后，发现仪表盘显示左前轮胎压为\,\qty{210}{kPa}，此时轮胎内气体温度为\,\qty{-13}{\degreeCelsius}，请判断轮胎是否漏气；如果漏气，求剩余气体与原来气体的质量之比。
\end{enumerate}